What was the last scientific article that you read? Think about it.
Not the last article that you paged through or
skimmed or sampled the introduction and conclusion, but one in which you
actually <i>read</i> most of the pages. Now, how long was that article? 
Was it a 10-12
conference article, a 20-page journal article, or a 50-page mini-monograph?

Glenn Ellison of the MIT Department of Economics presents compelling
evidence (<a href="http://econ-www.mit.edu/faculty/gellison/files/jrnth.pdf"><em>Evolving Standards for Academic Publishing: A q-r Theory</em></a>) that
"In almost all fields papers seem to be longer now than in 1975." Ellison
sampled journals in some thirty-two fields and found significant increases
in average page length. For example, ACM <em>JACM</em> papers increased in
length from 12 pages to 29 pages (though <em>CACM</em> papers actually
decreased from 7.6 pages on average to 7.0 pages).

Figure 1 provides more detail for ACM <em>TODS</em>. The top
line states the length in pages of the longest article in each yearly
volume, the middle line indicates the average length, and the bottom
line states the length of the shortest article.

<center>
<a name="FigureOne">
<img src="dec01.jpg" alt="Figure 1" width="100" border="0">
</a>
<br>
<b>Figure 1:colon[]</b> Article length per volume in ACM <em>TODS</em>
</center>

Quite frankly, all three trends are disturbing. The average article length
has more than doubled, from 19.2 pages in 1976 to 41.9 pages this year. The
average article this year was longer than the longest article in 1976. The
shortest article this year, at 31 pages, was longer than the average article
for the entire first decade of <em>TODS</em>' existence. In five separate years
an article of at least 60 pages appeared (one weighed in at a whopping 79
pages).

The total page count per volume was relatively stable over this period,
which meant that an increasing average article length was coupled with a
decreasing number of articles, as shown in <a href="#FigureTwo">Figure 2</a>. 
The average
issue of <em>TODS</em> in the eighties contained six articles, while in the
last seven years the journal has averaged only three articles per quarterly
issue, or a paper a month.

<center>
<a name="FigureTwo">
<img src="dec01.jpg" alt="Figure 2" width="100" border="0">
</a>
<br>
<b>Figure 2:colon[]</b> Number of articles per volume in ACM <em>TODS</em>
</center>

The result is that readers are confronted with less diverse and more
ponderous papers in each issue, of concern to the Editorial Board. (A
related phenomenon, also of concern, is that referees are confronted with
longer manuscripts, which must lengthen reviewing time.) Hence, the question
that opens this essay.

The cause of these trends is clear: in an effort to achieve
the "significance of contribution" required of acceptance, authors
are adding more detail, more theorems, more performance studies, and more
discussion to their submissions. (I
sheepishly admit that one of the data points in the top line in
<a href="#FigureOne">Figure 1</a>
is a paper of mine.)

It is less obvious that longer papers are better papers.
The top five most-cited <em>TODS</em> articles, by Chen, Smith&Smith,
Codd, Shipman, and Stonebraker/Wong/Kreps/Held, average a hair over 30 pages
each (again, less than the shortest article this year).

Over-simplifying, there are five basic types of papers,
appropriate for different venues. There are highly innovative, initial-cut
papers that have neat ideas that haven't yet been fleshed out; in fact, it
may turn out that the idea doesn't work. Workshops are an ideal venue for
such papers. The next step is to pursue the idea further, prove some things
or test the concepts out in an implementation, realizing a paper appropriate
for a conference.  An idea that has merit can be elaborated and detailed
evaluation undertaken, generating significant theorems and/or implementation
and/or analytical or empirical performance studies, realizing a paper
appropriate for a journal such as <em>TODS</em>. In fact, the key role of
archival journals is to publish solid, enduring, reproducible research in a
format less constraining than workshop or conference proceedings (and
benefitting from a more thorough, multi-pass review).

The following policy has been in place for many years, somewhat stemming
the growth of the longest papers.

:quote[
6. <em>TODS</em> will discourage excessively long papers (longer than 50
double-spaced pages including figures, references, etc.), and unnecessary
digressions even in shorter papers. This is to motivate the authors to bring
out the essence of their papers more clearly, to make it easier for the
reviewers and readers, and to allow <em>TODS</em> to publish more papers in
any given issue.
]

This policy states clearly the advantages of keeping journal papers focused,
while retaining their positive qualities of completeness and rigor.

The fourth kind of paper is a survey of the literature or of practice. My
previous editorial announced that <em>TODS</em> is now accepting certain
kinds of directed surveys.

The final kind of paper has traditionally appeared in journals such as 
<em>Information Processing Letters</em> or 
ACM <em>Letters on Programming Languages and Systems</em>. 
These papers are also rigorous and contribute e.g.&nbsp;a
faster algorithm or a theoretical insight that has not appeared
elsewhere, yet have a more restricted scope than a conference or
conventional journal paper. They are usually quite short, 5-20 pages,
depending on the background they need to present.

The Editorial Board recently augmented the above maximum bound with the
<a href="http://www.acm.org/tods/Authors.htm">following policy</a> 
which focuses on the minimum bound.

:quote[
7. In a similar vein, <em>TODS</em> encourages shorter submissions, including
even very short (say, five page) submissions. The primary criterion for
acceptance is improving on the state-of-the-art in some significant way.
]

We are serious about encouraging shorter submissions. In fact, we are
delighted to have received recently a submission of five pages and another
of seven pages. We will inform referees of this new policy and will strive
to reward succinct submissions with faster turnaround and favorable editorial
decisions. We hope that authors who had decided their submission wasn't
long enough for <em>TODS</em> will rethink that decision in light of this new
policy.

On a related note, you can now submit your manuscript electronically, in
a variety of formats. Just point your browser at the above-mentioned URL.
